\documentclass[11pt]{article}
\usepackage[letterpaper,margin=0.86in]{geometry}
\usepackage[T1]{fontenc}
\usepackage{lmodern,microtype}
\usepackage{amsmath,amssymb,amsthm,mathtools,bm}
\usepackage{booktabs,array,enumitem}
\usepackage[hidelinks]{hyperref}
\hypersetup{
 pdftitle={Exact Zero-Frequency Meissner Weight in a Gapless Frustration-Free QGN Parent},
 pdfauthor={Nicholas Sledgianowski},
 pdfsubject={Graph-Hodge response and zero-frequency Meissner weight}
}
\usepackage[nameinlink,capitalise]{cleveref}
\allowdisplaybreaks
\setlength{\emergencystretch}{2em}

\newtheorem{theorem}{Theorem}[section]
\newtheorem{proposition}[theorem]{Proposition}
\newtheorem{lemma}[theorem]{Lemma}
\newtheorem{corollary}[theorem]{Corollary}
\theoremstyle{definition}
\newtheorem{definition}[theorem]{Definition}
\theoremstyle{remark}
\newtheorem{remark}[theorem]{Remark}

\newcommand{\cH}{\mathcal H}
\newcommand{\cK}{\mathcal K}
\newcommand{\cC}{\mathcal C}
\newcommand{\cL}{\mathcal L}
\newcommand{\cZ}{\mathcal Z}
\newcommand{\cT}{\mathcal T}
\newcommand{\cS}{\mathcal S}
\newcommand{\Id}{\mathbb I}
\newcommand{\Tr}{\operatorname{Tr}}
\newcommand{\Ran}{\operatorname{Ran}}
\newcommand{\Ker}{\operatorname{Ker}}
\newcommand{\spec}{\operatorname{spec}}
\newcommand{\Span}{\operatorname{span}}
\newcommand{\dd}{\mathrm d}
\newcommand{\I}{\mathrm i}
\newcommand{\e}{\mathrm e}
\newcommand{\ket}[1]{\lvert #1\rangle}
\newcommand{\bra}[1]{\langle #1\rvert}
\newcommand{\abs}[1]{\lvert #1\rvert}
\newcommand{\norm}[1]{\lVert #1\rVert}
\newcommand{\ip}[2]{\langle #1,#2\rangle}
\newcommand{\braket}[1]{\langle #1\rangle}
\newcommand{\Pcyc}{P_{\mathrm{cyc}}}
\newcommand{\Pker}{P_{\ker}}

\title{\textbf{Exact Zero-Frequency Meissner Weight\\
in a Gapless Frustration-Free QGN Parent}\\[0.35em]
\large Graph-Hodge response, sharp Kohn--Drude limit criteria,\\
and protected multiblock stiffness}
\author{Nicholas Sledgianowski}
\date{Reviewer-revised research draft, July 27, 2026}

\begin{document}
\maketitle

\begin{abstract}
Finite-volume Peierls curvature is an adiabatic Kohn response, not automatically a bulk zero-frequency
Meissner or superfluid weight.  The thermodynamic, frequency, and transverse-momentum limits can fail
to commute.  We separate and solve these issues in a frustration-free paired lattice setting.

First, for an arbitrary analytic finite-volume ground cluster, the imaginary-frequency current kernel
exceeds the Kohn curvature by an explicit positive spectral defect.  After volume normalization, the
thermodynamic-first Abelian kernel equals the thermodynamic Kohn curvature if and only if the
current-generated spectral measure is uniformly tight at zero energy in the $H^{-1}$ norm, assuming a
uniform total $H^{-1}$ bound.  A uniform current-sector gap is sufficient but not necessary.  For every
positive-square parent $H(A)=D(A)^\dagger D(A)/2$, the entire finite-frequency response is the exact
regularized least-squares residual
\[
 \cK_L(\I\zeta)
 =S_L^\dagger\frac{\zeta^2}{\cL_L^2+\zeta^2}S_L,
 \qquad
 S_L=(\partial_A D_L)(0)P_L,
 \quad
 \cL_L=\frac12D_LD_L^\dagger.
\]

Second, we close the independent transverse limit for an explicit local paired parent.  On a connected
graph, complete gauge-covariant site-swap rows turn the static least-squares subtraction into graph Hodge
decomposition.  In the fixed-$n$ Dicke/AGP ground state,
\[
 \cC_{G,n}[a]
 =j\gamma_{N,n}\ip{a}{P_{\ker B}a},
 \qquad
 \cK_{G,n}(\I\zeta;a)
 =j\gamma_{N,n}\ip{a}{
 \frac{\zeta^2}{(jB^*B/4)^2+\zeta^2}a},
 \qquad
 \gamma_{N,n}=\frac{2n(N-n)}{N(N-1)},
\]
where $B$ is the oriented incidence matrix.  Every lattice-transverse field satisfies $Ba=0$, so its
static and dynamical kernels coincide at every finite size, wave vector, and positive frequency regulator.
On a periodic lattice in $d\ge2$, with $n/N\to\rho\in(0,1)$,
\[
 \boxed{\kappa_{\rm K}=D_{\rm A}=D_{\rm M}=2j\rho(1-\rho)>0}
\]
in the electronic-link convention, or $D_s=j\rho(1-\rho)/2$ in the pair-phase convention.  The same
state has exact pair off-diagonal long-range order, while its fixed-number gap closes as $L^{-2}$; no
uniform many-body gap is used.  Finally, any multiblock positive-square extension containing the swap
backbone inherits a frequency-independent positive transverse operator floor.  The result is a genuine
zero-temperature matter Meissner theorem for the exact parent.  Transfer to finite-coupling microscopic
downfoldings is deliberately left to a separate stability analysis.
\end{abstract}

\begin{quote}\small
\textbf{Claim hierarchy.}
\begin{enumerate}[leftmargin=1.6em,itemsep=0.12em,topsep=0.2em]
\item The finite-volume Kubo--Kohn identity, the positive defect, the $H^{-1}$ criterion, and the
positive-square dynamical least-squares formula are exact finite-dimensional theorems.
\item Under a uniform total current $H^{-1}$ bound, temporal Kohn--Drude equality is equivalent to
uniform low-energy $H^{-1}$ tightness.  A closing complete many-body gap is allowed.
\item A Meissner statement requires an independent transverse limit.  Temporal equality alone is not
enough; a clean normal metal is a counterexample.
\item For the gauge-covariant swap-QGN/AGP parent, both temporal and transverse problems are solved
algebraically at finite size.  The resulting Meissner weight is positive at every interior filling.
\item Multiblock extensions inherit a protected transverse floor in every selected product composition.
No scalar lowest-branch claim is made without a branch or state selector.
\item No finite-temperature theorem, unrestricted real-axis conductivity decomposition, material
parameter extraction, or finite-coupling microscopic transfer theorem is claimed here.
\end{enumerate}
\end{quote}

\section{Introduction}

Exact paired parent Hamiltonians make static response unusually accessible.  A complete Peierls twist
defines an analytic finite-volume Hamiltonian; frustration freeness identifies the exact zero manifold;
and second-order perturbation theory reduces the curvature to a positive least-squares residual.  These
ingredients underlie quantum-geometric-nesting (QGN) constructions, projected attractive-Hubbard
models, and recent lower bounds on flat-band stiffness \cite{Han2024,Gao2026}.  They also meet the
broader flat-band superfluid-weight literature, where quantum geometry supplies rigorous or mean-field
bounds rather than the exact many-body response identity obtained below
\cite{Peotta2015,Torma2022}.

The remaining distinction is physical rather than terminological.  Kohn curvature is obtained by taking
the frequency to zero at fixed finite volume and only then sending the volume to infinity
\cite{Kohn1964}.  A bulk uniform electric response takes the thermodynamic limit first.  A Meissner
response is a static transverse response at nonzero wave vector, followed by the transverse
zero-momentum limit.  The three limits diagnose different phenomena and are distinguished in the
standard current-correlation criteria for an insulator, metal, or superconductor
\cite{Scalapino1992,Scalapino1993,Resta2018,Watanabe2020}.  Tada and Koma's no-go analysis likewise
shows why the order of limits and the class of applied fields must be specified in a superconductivity
claim \cite{Tada2016}.  Explicit examples are known in which finite-size Kohn and thermodynamic Drude
orders differ \cite{Takasan2023}; even when they agree, a clean normal metal can have Drude weight but
no transverse Meissner response.

This paper gives a standalone resolution in two layers.

The first layer is general.  We derive a positive operator measuring the difference between the
finite-frequency kernel and static Kohn curvature.  Its thermodynamic limit vanishes precisely when the
current-generated low-energy spectral measure is uniformly integrable after weighting by inverse
excitation energy.  For positive-square parents, the same statement becomes a source-singular-mode
criterion in the target space of the null map.

The second layer is model specific and exact.  We introduce a gauge-covariant complete site-swap parent
on a connected graph.  Its seniority-zero compression is the symmetric-exclusion or ferromagnetic
hard-core-boson generator associated with the Matsubara--Matsuda correspondence
\cite{Matsubara1956}.  The current source is the graph divergence of the bond field.  Static relaxation is
therefore precisely Hodge projection onto gradients, while the residual is the graph cycle-space norm.
Every divergence-free profile lies in the target kernel and has a frequency-independent response.  On a
periodic lattice this includes all lattice-transverse Fourier modes and the uniform harmonic flux.

The outcome is an exact gapless example with
\[
 \kappa_{\rm K}=D_{\rm A}=D_{\rm M}>0.
\]
The paired Dicke state also has Yang off-diagonal long-range order \cite{Yang1962}, but the response
identity is proved independently of that correlator.  This distinction complements results deriving a
Meissner effect from ODLRO under additional algebraic or thermodynamic hypotheses
\cite{Sewell1990,Nieh1995}: here the finite-size current kernel itself is evaluated exactly before any
thermodynamic limit.  The state is also a zero-momentum relative of the $\eta$-pairing family, whose ODLRO
does not by itself settle the transport question \cite{Yang1989}.  The fixed-number spectral gap closes
quadratically; the theorem works because gauge geometry separates transverse sources from the soft
longitudinal sector, not because a hidden uniform gap survives.

\section{Three responses and their limit orders}
\label{sec:responses}

\subsection{Gauge-covariant finite tori and profile-dependent kernels}

Let $\Lambda_L$ be a periodic $d$-dimensional lattice of volume $V_L$, and fix a conserved
particle-number sector.  Let $H_L[A]$ be a finite-range or uniformly exponentially local Hamiltonian
coupled to a real bond vector potential by a complete Peierls prescription.  Gauge covariance means
\begin{equation}
 H_L[A+\nabla\phi]=U_\phi H_L[A]U_\phi^\dagger.
 \label{eq:gauge-covariance}
\end{equation}
For any real bond profile $a$ on $\Lambda_L$, set
\begin{equation}
 H_L(t;a)=H_L[ta],
 \qquad
 J_L[a]=\partial_tH_L(t;a)|_{t=0},
 \qquad
 T_L[a]=\partial_t^2H_L(t;a)|_{t=0}.
 \label{eq:JT-profile}
\end{equation}
A uniform flat connection in direction $v$ is the harmonic profile
$a^{\rm harm}_{L,v}$; when no argument is displayed below,
$J_L$, $T_L$, $\cC_L$, and $\cK_L$ refer to that profile.

Let $E_L$ be the lowest zero-field energy, $P_L$ the projector onto the complete finite-volume ground
cluster, $Q_L=\Id-P_L$, and
\begin{equation}
 \widehat H_L=H_L[0]-E_L\succeq0.
 \label{eq:Hhat}
\end{equation}
The cluster is assumed isolated for each finite $L$, but no size-independent gap is assumed.  For each
profile used in a response coefficient we impose
\begin{equation}
 P_LJ_L[a]P_L=0.
 \label{eq:no-linear}
\end{equation}
This holds automatically for positive-square parents and for time-reversal-symmetric zero-current
branches.  If it fails, first-order splitting or a cusp can precede any quadratic response.

\subsection{Kohn, Abelian, and transverse Meissner kernels}

For an arbitrary profile $a$, finite-dimensional degenerate perturbation theory gives the second-order
Kohn operator
\begin{equation}
 \cC_L[a]
 =P_LT_L[a]P_L-2P_LJ_L[a]Q_L\widehat H_L^+Q_LJ_L[a]P_L,
 \label{eq:Kohn-operator}
\end{equation}
and, for $\zeta>0$, the imaginary-frequency or Abel-regularized kernel
\begin{equation}
 \cK_L(\I\zeta;a)
 =P_LT_L[a]P_L
 -2P_LJ_L[a]Q_L\frac{\widehat H_L}{\widehat H_L^2+\zeta^2}Q_LJ_L[a]P_L.
 \label{eq:Abelian-kernel}
\end{equation}
Write
\begin{equation}
 c_L[a]=\frac{\cC_L[a]}{V_L},
 \qquad
 k_L(\zeta;a)=\frac{\cK_L(\I\zeta;a)}{V_L},
 \qquad
 k_L(0;a):=c_L[a].
 \label{eq:profile-normalized}
\end{equation}
The value at $\zeta=0$ is therefore, by definition, the finite-volume static curvature; it is not a
thermodynamic-first frequency limit.

For the harmonic uniform profile, abbreviate
$c_L(v)=c_L[a^{\rm harm}_{L,v}]$ and
$k_L(\zeta;v)=k_L(\zeta;a^{\rm harm}_{L,v})$.  When the limits exist, define
\begin{align}
 \kappa_{\rm K}(v)&=\lim_{L\to\infty}c_L(v),
 \label{eq:kappaK}\\
 D_{\rm A}(v)&=\lim_{\zeta\downarrow0}\lim_{L\to\infty}k_L(\zeta;v).
 \label{eq:DA}
\end{align}
The first takes the static limit at finite volume; the second takes the thermodynamic limit first.

For $d\ge2$, let $a_{L,q,v}$ be a real lattice-transverse bond profile with Fourier support at $\pm q$,
polarization $v$, exact lattice divergence zero, and normalization
$\norm{a_{L,q,v}}^2=V_L$.  Define
\begin{equation}
 \cK_{L,T}(q,\I\zeta;v):=\cK_L(\I\zeta;a_{L,q,v}),
 \qquad
 k_{L,T}(q,\zeta;v):=\frac{\cK_{L,T}(q,\I\zeta;v)}{V_L},
 \label{eq:finite-q-kernel}
\end{equation}
with
\begin{equation}
 \cK_{L,T}(q,0;v):=\cC_L[a_{L,q,v}],
 \qquad
 k_{L,T}(q,0;v):=c_L[a_{L,q,v}].
 \label{eq:finite-q-static-convention}
\end{equation}
Thus every occurrence of $k_{L,T}(q,0;v)$ means that $\zeta\downarrow0$ has already been taken at
fixed finite $L$.

The allowed momenta depend on $L$.  An admissible transverse zero-momentum limit is a sequence
$q_m\to0$ together with, for each $m$, a cofinal subsequence of tori on which
$q_m\in\Lambda_L^*$ and the profiles $a_{L,q_m,v}$ are defined.  For rational reciprocal vectors one
may, for example, take $L$ through integer multiples of a common period.  The matter Meissner weight is
\begin{equation}
 D_{\rm M}(v)
 =\lim_{m\to\infty}\lim_{\substack{L\to\infty\\q_m\in\Lambda_L^*}}
 k_{L,T}(q_m,0;v),
 \label{eq:DM}
\end{equation}
when the result is independent of the admissible sequence.  The alternative quantity
\begin{equation}
 \lim_{m\to\infty}\lim_{\zeta\downarrow0}\lim_{\substack{L\to\infty\\q_m\in\Lambda_L^*}}
 k_{L,T}(q_m,\zeta;v)
\end{equation}
is a separate order of limits unless temporal uniformity has also been established at nonzero $q$.  In the
exact swap model below the two definitions coincide already at finite size because the transverse kernel is
independent of $q$ and $\zeta$.

The adjective ``transverse'' is lattice exact: the polarization is annihilated by the lattice divergence
symbol, which approaches $q\cdot v=0$ at small momentum.

\begin{remark}[Matter response]
The coefficient in \eqref{eq:DM} is the zero-temperature matter current kernel before coupling to the
Maxwell field.  Conventional factors of electric charge and optical $\pi$ multiply all three coefficients
and do not affect the order-of-limits statements.
\end{remark}

\section{Exact finite-volume Kubo--Kohn identity}
\label{sec:finite-identity}

\begin{theorem}[Positive finite-frequency defect]
\label{thm:positive-defect}
Under the hypotheses of \cref{sec:responses}, for every finite $L$, every profile $a$ satisfying
\eqref{eq:no-linear}, and every $\zeta>0$,
\begin{equation}
 \boxed{
 \cK_L(\I\zeta;a)-\cC_L[a]
 =2\zeta^2P_LJ_L[a]Q_L\widehat H_L^+
 (\widehat H_L^2+\zeta^2)^{-1}Q_LJ_L[a]P_L\succeq0.}
 \label{eq:positive-defect}
\end{equation}
Consequently $\cK_L(\I\zeta;a)$ decreases in the Loewner order to $\cC_L[a]$ as
$\zeta\downarrow0$.
\end{theorem}

\begin{proof}
Subtract \eqref{eq:Kohn-operator} from \eqref{eq:Abelian-kernel}.  On the positive spectrum,
\begin{equation}
 \frac1E-\frac{E}{E^2+\zeta^2}
 =\frac{\zeta^2}{E(E^2+\zeta^2)}\ge0.
\end{equation}
Functional calculus gives \eqref{eq:positive-defect}.  The scalar multiplier is increasing in $\zeta$,
which gives monotonicity.
\end{proof}

The theorem makes the finite-size statement exact.  At fixed $L$, static curvature and the
zero-frequency Kubo kernel coincide.  The only temporal issue is uniformity in $L$.

\subsection{Operator-valued current spectral measure}

Let $P_{\widehat H_L}(B)$ be the spectral projection of $\widehat H_L$ for a Borel set
$B\subset(0,\infty)$.  Define the positive operator-valued measure on $P_L\cH_L$ by
\begin{equation}
 \Sigma_L(B)
 =\frac1{V_L}P_LJ_LQ_LP_{\widehat H_L}(B)Q_LJ_LP_L.
 \label{eq:spectral-measure}
\end{equation}
Also put $d_L=V_L^{-1}P_LT_LP_L$.  Then
\begin{align}
 c_L&=d_L-2\int_{(0,\infty)}\frac1E\,\dd\Sigma_L(E),
 \label{eq:c-spectral}\\
 k_L(\zeta)&=d_L-2\int_{(0,\infty)}\frac{E}{E^2+\zeta^2}\,\dd\Sigma_L(E),
 \label{eq:k-spectral}\\
 k_L(\zeta)-c_L
 &=2\int_{(0,\infty)}\frac{\zeta^2}{E(E^2+\zeta^2)}\,\dd\Sigma_L(E).
 \label{eq:defect-spectral}
\end{align}
At finite volume these are finite sums.  State-resolved formulas follow by taking expectations in a
selected ground-state density matrix.

\begin{lemma}[The relevant gap is a current-sector gap]
\label{lem:charge-gap}
Suppose $[H_L,N_L]=0$ and the Peierls coupling preserves particle number.  Then every matrix element in
\eqref{eq:spectral-measure} lies in the same fixed-number sector as the ground cluster.  An
addition/removal charge gap therefore does not bound the denominators in \eqref{eq:c-spectral}.  A
uniform gap in the current-generated fixed-number subspace is sufficient.
\end{lemma}

\begin{proof}
Particle-number conservation gives $[J_L,N_L]=0$.  Hence $J_LP_L$ lies in the same number sector as
$P_L$, and only that sector enters the spectral measure.
\end{proof}

\section{Sharp temporal order-of-limits criterion}
\label{sec:Hminus1}

\begin{definition}[$H^{-1}$ boundedness and tightness]
Define
\begin{equation}
 M=\sup_L\left\|\int_{(0,\infty)}E^{-1}\,\dd\Sigma_L(E)\right\|.
 \label{eq:Hminus1-bound}
\end{equation}
The current measures are uniformly $H^{-1}$-tight at zero (equivalently, uniformly $H^{-1}$ equi-integrable at zero energy) if
\begin{equation}
 \lim_{\delta\downarrow0}\limsup_{L\to\infty}
 \left\|\int_{(0,\delta]}E^{-1}\,\dd\Sigma_L(E)\right\|=0.
 \label{eq:Hminus1-tightness}
\end{equation}
\end{definition}

\begin{theorem}[Sharp temporal criterion]
\label{thm:sharp-temporal}
Assume $M<\infty$.  Then the following are equivalent:
\begin{enumerate}[label=\textup{(\roman*)},leftmargin=2.6em]
\item the current measures satisfy \eqref{eq:Hminus1-tightness};
\item the finite-frequency defect vanishes uniformly in the thermodynamic limit,
\begin{equation}
 \lim_{\zeta\downarrow0}\limsup_{L\to\infty}
 \norm{k_L(\zeta)-c_L}=0.
 \label{eq:uniform-defect-zero}
\end{equation}
\end{enumerate}
More quantitatively, put
\begin{equation}
 t(\delta)=\limsup_{L\to\infty}
 \left\|\int_{(0,\delta]}E^{-1}\,\dd\Sigma_L(E)\right\|.
\end{equation}
Then
\begin{equation}
 t(\zeta)
 \le\limsup_{L\to\infty}\norm{k_L(\zeta)-c_L},
 \label{eq:lower-tail-bound}
\end{equation}
and, for $0<\zeta<\delta$,
\begin{equation}
 \limsup_{L\to\infty}\norm{k_L(\zeta)-c_L}
 \le2t(\delta)+2\frac{\zeta^2}{\delta^2}M.
 \label{eq:upper-tail-bound}
\end{equation}
\end{theorem}

\begin{proof}
For $0<E\le\zeta$,
\begin{equation}
 2\frac{\zeta^2}{E(E^2+\zeta^2)}\ge\frac1E,
\end{equation}
so positivity of the operator-valued measure gives \eqref{eq:lower-tail-bound}.  To prove
\eqref{eq:upper-tail-bound}, split \eqref{eq:defect-spectral} at $\delta$.  Below $\delta$, the scalar
kernel is at most $2/E$.  Above $\delta$, it is at most
$2(\zeta^2/\delta^2)E^{-1}$.  The triangle inequality for positive operator integrals yields the bound.
The equivalence follows by first taking $L\to\infty$, then $\zeta\downarrow0$, and in the upper bound
finally $\delta\downarrow0$.
\end{proof}

\begin{corollary}[Kohn equals the bulk Abelian kernel]
\label{cor:Kohn-Drude}
Assume \cref{thm:sharp-temporal}, assume $c_L\to\kappa_{\rm K}$ in operator norm, and assume
$k_L(\zeta)$ has a thermodynamic limit for every fixed $\zeta>0$.  Then
\begin{equation}
 \boxed{D_{\rm A}=\kappa_{\rm K}.}
\end{equation}
The same statement holds after taking an expectation in any consistently selected sequence of states or
density matrices in the ground cluster.
\end{corollary}

\begin{corollary}[Uniform current-sector gap]
\label{cor:current-gap}
If the spectral support of $Q_LJ_LP_L$ lies in $[\Delta,\infty)$ with $\Delta>0$ independent of $L$,
and if \eqref{eq:Hminus1-bound} holds, then \eqref{eq:Hminus1-tightness} holds.  The converse is false:
the full gap may close while the current matrix elements vanish sufficiently rapidly.
\end{corollary}

\begin{remark}[Mode-by-mode diagnostic]
For a selected scalar ground state, let a low mode have energy $\Delta_L$ and normalized current weight
\begin{equation}
 w_L=\frac1{V_L}\abs{\bra{m_L}J_L\ket{\Omega_L}}^2.
\end{equation}
Its contribution to the tightness functional is $w_L/\Delta_L$.  The mode is harmless if this ratio tends
to zero and obstructing if it has a positive limit.  The theorem is therefore a matrix-element statement,
not merely a dispersion statement.
\end{remark}

\section{Positive-square parents: dynamical least squares}
\label{sec:least-squares}

Let
\begin{equation}
 H_L(a)=\frac12D_L(a)^\dagger D_L(a),
 \qquad
 D_L(a):\cH_L\to\cT_L
 \label{eq:positive-square}
\end{equation}
be analytic, and let $P_L=P_{\ker D_L(0)}$.  Positivity gives $E_L=0$, so throughout this section
$\widehat H_L=H_L(0)$.  Put
\begin{equation}
 D_L=D_L(0),
 \qquad
 S_L=(\partial_aD_L)(0)P_L,
 \qquad
 \cL_L=\frac12D_LD_L^\dagger.
 \label{eq:source-target-Laplacian}
\end{equation}
The map $S_L$ is the stacked Peierls source and $\cL_L$ is the target-space Laplacian.

\begin{theorem}[Dynamical least-squares identity]
\label{thm:dynamical-LS}
For every finite $L$ and $\zeta>0$,
\begin{equation}
 \boxed{
 \cK_L(\I\zeta)
 =S_L^\dagger\frac{\zeta^2}{\cL_L^2+\zeta^2}S_L.}
 \label{eq:dynamical-LS}
\end{equation}
Its static Kohn curvature is
\begin{equation}
 \boxed{\cC_L=S_L^\dagger P_{\ker D_L^\dagger}S_L.}
 \label{eq:static-LS}
\end{equation}
Equivalently, for $\psi\in P_L\cH_L$,
\begin{equation}
 \bra{\psi}\cK_L(\I\zeta)\ket{\psi}
 =\min_{y\in\cT_L}
 \left\{\norm{S_L\psi-\cL_Ly}^2+\zeta^2\norm{y}^2\right\}.
 \label{eq:Tikhonov}
\end{equation}
\end{theorem}

\begin{proof}
Since $D_LP_L=0$,
\begin{equation}
 J_LP_L=\frac12D_L^\dagger S_L,
 \qquad
 P_LT_LP_L=S_L^\dagger S_L.
\end{equation}
Singular-value functional calculus gives
\begin{equation}
 \frac12D_L\frac{\widehat H_L}{\widehat H_L^2+\zeta^2}D_L^\dagger
 =\frac{\cL_L^2}{\cL_L^2+\zeta^2}.
\end{equation}
Substitution into \eqref{eq:Abelian-kernel} proves \eqref{eq:dynamical-LS}.  Taking $\zeta\downarrow0$
projects onto $\ker\cL_L=\ker D_L^\dagger$ and gives \eqref{eq:static-LS}.  Minimizing the quadratic form
in \eqref{eq:Tikhonov} gives the normal equation
$(\cL_L^2+\zeta^2)y=\cL_LS_L\psi$ and reproduces \eqref{eq:dynamical-LS}.
\end{proof}

\begin{corollary}[Sharp source-singular-mode criterion]
\label{cor:source-criterion}
Assume
\begin{equation}
 \sup_L\frac1{V_L}\norm{S_L^\dagger S_L}<\infty.
\end{equation}
Then temporal Kohn--Drude equality holds if and only if
\begin{equation}
 \lim_{\delta\downarrow0}\limsup_{L\to\infty}
 \frac1{V_L}\norm{S_L^\dagger\mathbf1_{(0,\delta]}(\cL_L)S_L}=0.
 \label{eq:source-tightness}
\end{equation}
Thus a closing positive singular-value sector is harmless precisely when it carries vanishing extensive
Peierls-source weight.
\end{corollary}

\begin{proof}
On the positive singular spectrum, $D_L$ identifies the state-space current measure with one half of the
target-space source measure.  Equivalently, subtract \eqref{eq:static-LS} from
\eqref{eq:dynamical-LS}; the defect multiplier is $\zeta^2/(\lambda^2+\zeta^2)$ and the same cutoff
argument as in \cref{thm:sharp-temporal} applies.
\end{proof}

\section{From a uniform kernel to a transverse Meissner weight}
\label{sec:transverse-general}

\begin{definition}[Static transverse continuity]
Choose a sequence of ground states or density matrices defining the thermodynamic phase.  The static
transverse response is continuous at zero momentum in direction $v$ if, for every admissible sequence
specified after \eqref{eq:finite-q-static-convention},
\begin{equation}
 \lim_{m\to\infty}\lim_{\substack{L\to\infty\\q_m\in\Lambda_L^*}}
 k_{L,T}(q_m,0;v)
 =\lim_{L\to\infty}c_L(v).
 \label{eq:transverse-continuity}
\end{equation}
\end{definition}

\begin{theorem}[Kohn--Drude--Meissner completion]
\label{thm:general-completion}
Suppose, in a selected thermodynamic state,
\begin{enumerate}[label=\textup{(\roman*)},leftmargin=2.6em]
\item the hypotheses of \cref{cor:Kohn-Drude} hold for the uniform current;
\item static transverse continuity \eqref{eq:transverse-continuity} holds.
\end{enumerate}
Then
\begin{equation}
 \boxed{\kappa_{\rm K}=D_{\rm A}=D_{\rm M}.}
 \label{eq:three-equal-general}
\end{equation}
\end{theorem}

\begin{proof}
The first equality is \cref{cor:Kohn-Drude}; the second is
\eqref{eq:transverse-continuity} and the definition \eqref{eq:DM}.
\end{proof}

The theorem is modular.  Temporal failure is low-energy inverse-energy current weight.  Spatial failure is
a discontinuity between an exactly uniform connection and the limit of nonzero transverse fields.  The two
obstructions are independent.

\begin{proposition}[A sufficient summability condition]
\label{prop:summability}
Let $j_v(x)$ be a local current density in a translation-invariant selected pure ground state and write
$J_{L,v}=\sum_xj_v(x)$.  Define the connected Euclidean correlation by
\begin{equation}
 G_{L,v}(x,t)
 =\bra{\Omega_L}j_v(x)\e^{-t\widehat H_L}Q_Lj_v(0)\ket{\Omega_L},
 \qquad t\ge0.
 \label{eq:Euclidean-correlation}
\end{equation}
The insertion of $Q_L$ is the connected subtraction.  Assume local thermodynamic limits exist,
translation invariance holds along the chosen sequence, and there is an $L$-independent nonnegative
envelope $F$ such that
\begin{equation}
 \abs{G_{L,v}(x,t)}\le F(x,t),
 \qquad
 \sum_x\int_0^\infty F(x,t)\,\dd t<\infty.
 \label{eq:summability}
\end{equation}
Assume also that the complete finite-$q$ diamagnetic term has a thermodynamic limit continuous at
$q=0$.  Then the uniform current is $H^{-1}$-tight and the static transverse kernel is continuous at
$q=0$.  The same conclusion holds for a finite ground multiplet if the bound is uniform for all matrix
elements in the selected density matrix.
\end{proposition}

\begin{proof}
Let $\mu_{L,v}$ be the scalar volume-normalized spectral measure of $J_{L,v}$ in the selected state.
Translation invariance and the spectral theorem give
\begin{equation}
 g_{L,v}(t)
 :=\frac1{V_L}\bra{\Omega_L}J_{L,v}\e^{-t\widehat H_L}Q_LJ_{L,v}\ket{\Omega_L}
 =\sum_xG_{L,v}(x,t)
 =\int_{(0,\infty)}\e^{-tE}\,\dd\mu_{L,v}(E).
 \label{eq:time-spectral}
\end{equation}
The connected subtraction is precisely the projector $Q_L$, so no zero-energy atom occurs.  The
summable envelope implies
$\abs{g_{L,v}(t)}\le\sum_xF(x,t)$ with an $L$-independent integrable majorant.
Moreover, for $0<E\le\delta$,
\begin{equation}
 \frac1E
 \le \mathrm e\int_{1/\delta}^{\infty}\e^{-tE}\,\dd t.
\end{equation}
Therefore positivity and \eqref{eq:time-spectral} imply
\begin{equation}
 \int_{(0,\delta]}\frac{\dd\mu_{L,v}(E)}E
 \le \mathrm e\int_{1/\delta}^{\infty}g_{L,v}(t)\,\dd t
 \le \mathrm e\sum_x\int_{1/\delta}^{\infty}F(x,t)\,\dd t.
 \label{eq:summability-tail}
\end{equation}
The right side tends to zero uniformly in $L$, proving the $H^{-1}$ tightness criterion.

For the static nonzero-momentum paramagnetic susceptibility, the same spectral representation gives
\begin{equation}
 \chi_{L,v}(q)
 =\sum_x\e^{-\I q\cdot x}\int_0^\infty G_{L,v}(x,t)\,\dd t.
 \label{eq:Fourier-susceptibility}
\end{equation}
Local thermodynamic convergence and \eqref{eq:summability} permit the $L\to\infty$ limit term by term;
absolute summability then permits $q\to0$ by dominated convergence.  Hence the thermodynamic
paramagnetic term is continuous at zero momentum.  Adding the assumed continuous diamagnetic term
proves \eqref{eq:transverse-continuity}.
\end{proof}

The exact model below uses a stronger mechanism: a finite-size source-kernel identity removes the need
for asymptotic correlation estimates.

\section{Gauge-covariant swap-QGN parent on a graph}
\label{sec:graph-model}

\subsection{Fermionic definition}

Let $G=(V,E)$ be a finite connected graph with $N=\abs V$ vertices and a fixed orientation of each edge.
At site $x$ place spinful fermions $c_{x\sigma}$, $\sigma=\uparrow,\downarrow$, with
\begin{equation}
 N_x=n_{x\uparrow}+n_{x\downarrow},
 \qquad
 p_x^\dagger=c_{x\uparrow}^\dagger c_{x\downarrow}^\dagger,
 \qquad
 q_x=(n_{x\uparrow}-n_{x\downarrow})^2.
\end{equation}
Let $W_e$ be the complete fermionic site swap on the endpoints of $e=(x\to y)$, so that it exchanges all
local creation operators at $x$ and $y$ and leaves the other sites fixed.  Define the Peierls-covariant swap
and projector row
\begin{equation}
 W_e[A]
 =\e^{-\I A_eN_y}W_e\e^{\I A_eN_y},
 \qquad
 \Pi_e[A]=\frac{\Id-W_e[A]}2.
 \label{eq:twisted-swap}
\end{equation}
Since $W_e[A]$ is a Hermitian involution, $\Pi_e[A]$ is a projector.  The parent is
\begin{equation}
 \boxed{
 H_G[A]
 =\frac U2\sum_{x\in V}q_x^2
 +\frac j2\sum_{e\in E}\Pi_e[A]^2,
 \qquad U,j>0.}
 \label{eq:swap-parent}
\end{equation}
It is a local positive-square Hamiltonian.

Let $B:\mathbb R^E\to\mathbb R^V$ be the oriented incidence matrix,
\begin{equation}
 B_{z e}=\begin{cases}
 -1,&z=t(e),\\
 +1,&z=h(e),\\
 0,&\text{otherwise}.
 \end{cases}
 \label{eq:incidence}
\end{equation}
For a site function $\phi$, $B^*\phi$ is the discrete gradient.

\begin{proposition}[Complete gauge covariance]
\label{prop:gauge-covariance}
With
\begin{equation}
 U_\phi=\exp\!\left(-\I\sum_x\phi_xN_x\right),
\end{equation}
one has
\begin{equation}
 U_\phi H_G[A]U_\phi^\dagger=H_G[A+B^*\phi].
 \label{eq:graph-gauge}
\end{equation}
\end{proposition}

\begin{proof}
For $e=(x\to y)$, conjugation by $U_\phi$ multiplies a fermion transferred from $x$ to $y$ by
$\e^{-\I(\phi_y-\phi_x)}$ and the reverse transfer by the conjugate phase.  This is exactly the replacement
$A_e\mapsto A_e+(B^*\phi)_e$ in \eqref{eq:twisted-swap}.  The onsite rows are gauge invariant.
\end{proof}

\subsection{Exact fixed-number ground state}

For $0\le n\le N$, define the normalized paired Dicke/AGP state
\begin{equation}
 \ket{\Omega_{N,n}}
 =\binom Nn^{-1/2}
 \sum_{\substack{S\subset V\\\abs S=n}}
 \prod_{x\in S}p_x^\dagger\ket0.
 \label{eq:Dicke}
\end{equation}

\begin{theorem}[Complete fixed-$2n$ zero space]
\label{thm:kernel}
At $A=0$, in the sector of total particle number $2n$,
\begin{equation}
 \ker H_G[0]\big|_{2n}=\Span\{\ket{\Omega_{N,n}}\}.
 \label{eq:kernel}
\end{equation}
Odd-particle sectors have no zero vectors.
\end{theorem}

\begin{proof}
The rows $q_x$ force each site to be empty or doubly occupied.  In that seniority-zero subspace, each
$W_e$ is the ordinary transposition of the hard-core pair occupations at the endpoints.  The rows
$\Pi_e$ require invariance under every edge transposition.  Because $G$ is connected, the edge
transpositions generate the full symmetric group on $V$, whose fixed subspace at pair number $n$ is the
one-dimensional uniform subset state \eqref{eq:Dicke}.  A state with odd total particle number contains
a singly occupied site and is penalized by an onsite row.
\end{proof}

\subsection{Seniority-zero compression and one-body functions}

Write a hard-core pair configuration as $\ket S$, $S\subset V$, $\abs S=n$.  In the paired sector,
\begin{equation}
 H_{G,n}
 =\frac j4\sum_{e\in E}(\Id-W_e).
 \label{eq:exclusion-generator}
\end{equation}
With $m_x=S_x^z+\tfrac12$ and $p_x^\dagger=S_x^+$, this is exactly the isotropic spin-$\tfrac12$
ferromagnet
\begin{equation}
 H_{G,n}=\frac j2\sum_{e=(x,y)}\left(\frac14-\bm S_x\!\cdot\!\bm S_y\right),
 \label{eq:ferromagnet-form}
\end{equation}
restricted to fixed magnetization.  The Peierls field twists the $XY$ exchange with charge two.  This
form makes both the quadratic magnon gap and the graph-Hodge separation of longitudinal and transverse
sources transparent.  The onsite coefficient $U>0$ is used only to select seniority zero; it drops out of
all response formulas inside the paired ground sector.
For a real site function $f$ with $\sum_x f_x=0$, define
\begin{equation}
 \ket f
 =\binom Nn^{-1/2}
 \sum_{\abs S=n}\left(\sum_{x\in S}f_x\right)\ket S.
 \label{eq:site-function-state}
\end{equation}
Let
\begin{equation}
 \alpha_{N,n}=\frac{n(N-n)}{N(N-1)},
 \qquad
 \gamma_{N,n}=2\alpha_{N,n}.
 \label{eq:alpha-gamma}
\end{equation}

\begin{lemma}[Exact site-function representation]
\label{lem:site-functions}
For zero-mean $f,g$,
\begin{align}
 \braket{f|g}&=\alpha_{N,n}\ip{f}{g},
 \label{eq:site-inner}\\
 H_{G,n}\ket f&=\frac j4\ket{BB^*f}.
 \label{eq:site-Laplacian}
\end{align}
\end{lemma}

\begin{proof}
For \eqref{eq:site-inner}, the probability that a uniformly chosen $n$-subset contains $x$ is $n/N$,
and the probability that it contains distinct $x,y$ is $n(n-1)/[N(N-1)]$.  Expanding the two subset
sums and using $\sum_x f_x=\sum_xg_x=0$ leaves the coefficient $\alpha_{N,n}$.

For \eqref{eq:site-Laplacian}, an edge transposition changes
$\sum_{z\in S}f_z$ by the endpoint difference whenever exactly one endpoint is occupied.  Summing over
edges gives the graph Laplacian $BB^*$ acting on $f$.
\end{proof}

\section{Exact graph-Hodge response}
\label{sec:Hodge-response}

Let $a\in\mathbb R^E$ be a real bond profile and set $A=ta$.  Define
\begin{equation}
 J_G[a]=\partial_tH_G[ta]|_{t=0},
 \qquad
 T_G[a]=\partial_t^2H_G[ta]|_{t=0}.
\end{equation}
The graph Hodge projector onto the cycle space is
\begin{equation}
 \Pcyc
 =P_{\ker B}
 =\Id-B^*(BB^*)^+B.
 \label{eq:Hodge-projector}
\end{equation}

\begin{lemma}[Exact current and stress sources]
\label{lem:source-stress}
On the fixed-$n$ ground state,
\begin{align}
 J_G[a]\ket{\Omega_{N,n}}
 &=\frac{\I j}{2}\ket{Ba},
 \label{eq:current-source}\\
 \bra{\Omega_{N,n}}T_G[a]\ket{\Omega_{N,n}}
 &=j\gamma_{N,n}\norm a^2.
 \label{eq:stress-expectation}
\end{align}
\end{lemma}

\begin{proof}
On an edge $e=(x\to y)$, the twisted swap multiplies a pair moved from $x$ to $y$ by
$\e^{-2\I t a_e}$ and the reverse move by $\e^{2\I t a_e}$.  Differentiating the uniform subset state
therefore produces the signed endpoint occupation difference.  Summing over edges gives the divergence
$Ba$ and the coefficient in \eqref{eq:current-source}.  A second derivative contributes only when the
edge has exactly one occupied endpoint.  That probability is
$2n(N-n)/[N(N-1)]=\gamma_{N,n}$, proving \eqref{eq:stress-expectation}.
\end{proof}

\begin{theorem}[Exact graph-Hodge static and dynamical kernels]
\label{thm:graph-Hodge}
For every finite connected graph, every pair number $0\le n\le N$, every real bond profile $a$, and every
$\zeta>0$,
\begin{align}
 \boxed{
 \cC_{G,n}[a]
 =j\gamma_{N,n}\ip{a}{\Pcyc a},}
 \label{eq:static-Hodge}\\
 \boxed{
 \cK_{G,n}(\I\zeta;a)
 =j\gamma_{N,n}
 \ip{a}{\frac{\zeta^2}{(jB^*B/4)^2+\zeta^2}a}.}
 \label{eq:dynamical-Hodge}
\end{align}
Consequently:
\begin{enumerate}[label=\textup{(\roman*)},leftmargin=2.6em]
\item if $a=B^*\phi$ is a gauge gradient, then $\cC_{G,n}[a]=0$;
\item if $Ba=0$, then
\begin{equation}
 \cK_{G,n}(\I\zeta;a)=\cC_{G,n}[a]
 =j\gamma_{N,n}\norm a^2
 \label{eq:transverse-exact}
\end{equation}
for every $\zeta>0$.
\end{enumerate}
\end{theorem}

\begin{proof}
By \cref{lem:source-stress}, the current source lies in the site-function sector.  On that sector,
\cref{lem:site-functions} identifies the Hamiltonian with $(j/4)BB^*$.  Hence the paramagnetic term in
\eqref{eq:Abelian-kernel} is
\begin{equation}
 \frac{j^2}{2}\alpha_{N,n}
 \ip{Ba}{
 \frac{(j/4)BB^*}{[(j/4)BB^*]^2+\zeta^2}Ba}.
\end{equation}
Using the common nonzero spectrum of $BB^*$ and $B^*B$, and subtracting this term from
\eqref{eq:stress-expectation}, gives \eqref{eq:dynamical-Hodge}.  Taking $\zeta\downarrow0$ projects
onto $\ker(B^*B)=\ker B$ and proves \eqref{eq:static-Hodge}.  The two consequences follow from the
orthogonal Hodge decomposition
$\mathbb R^E=\Ran B^*\oplus\ker B$.
\end{proof}

\begin{remark}[Exact finite-frequency graph filter]
Equation \eqref{eq:dynamical-Hodge} is the model's central infrared mechanism.  On an edge-Laplacian
eigenmode with eigenvalue $\lambda>0$, the dynamical multiplier is
\begin{equation}
 \frac{\zeta^2}{(j\lambda/4)^2+\zeta^2}.
\end{equation}
At fixed $\zeta$, sufficiently soft longitudinal modes remain unresolved; the static limit removes them.
Cycle-space modes have $\lambda=0$ and survive with multiplier one at every frequency.  The general
$H^{-1}$ criterion is therefore realized here as an exact graph filter, without a uniform many-body gap.
For a divergence-free field, \eqref{eq:current-source} gives
$J_G[a]\ket{\Omega_{N,n}}=0$ already at finite size.  The transverse response is consequently purely
diamagnetic, independent of $q$, and independent of $\zeta$.  This is a deliberately fine-tuned feature
of the frustration-free parent, stronger than the partial diamagnetic--paramagnetic cancellation in a
generic BCS state.
\end{remark}

\section{Zero-frequency Meissner theorem on periodic lattices}
\label{sec:periodic-Meissner}

Let $G_L$ be the nearest-neighbor graph of a periodic $d$-dimensional Bravais lattice with
$N=L^d$ sites.  A Fourier bond field $a_\mu(q)$ is lattice transverse when
\begin{equation}
 \sum_{\mu=1}^d(1-\e^{-\I q_\mu})a_\mu(q)=0.
 \label{eq:lattice-transverse}
\end{equation}
For every allowed nonzero $q$ in $d\ge2$, there is a nonzero transverse polarization.  Normalize the
corresponding real profile by $\norm a^2=N$.

\begin{theorem}[Exact zero-frequency Meissner weight]
\label{thm:exact-Meissner}
Let $d\ge2$, let $n_L/N\to\rho\in(0,1)$, and consider the parent \eqref{eq:swap-parent} on $G_L$.
Then every finite-volume uniform harmonic profile, every allowed lattice-transverse Fourier profile, and
every $\zeta>0$ have the same response density,
\begin{equation}
 \frac1N\cK_{G_L,n_L}(\I\zeta;a)
 =\frac1N\cC_{G_L,n_L}[a]
 =j\gamma_{N,n_L}.
 \label{eq:finite-Meissner-density}
\end{equation}
All admissible orders of $L\to\infty$, $\zeta\downarrow0$, and transverse $q\to0$ therefore agree:
\begin{equation}
 \boxed{
 \kappa_{\rm K}=D_{\rm A}=D_{\rm M}=2j\rho(1-\rho)>0.}
 \label{eq:exact-Meissner-value}
\end{equation}
If the pair phase is $Q=2A$, the corresponding pair-phase stiffness is
\begin{equation}
 \boxed{D_s=\frac j2\rho(1-\rho).}
 \label{eq:pair-phase-value}
\end{equation}
\end{theorem}

\begin{proof}
A constant directional bond profile on a torus is divergence free, as is every field satisfying
\eqref{eq:lattice-transverse}.  Equation \eqref{eq:transverse-exact} and the normalization
$\norm a^2=N$ give \eqref{eq:finite-Meissner-density}.  Since
\begin{equation}
 \gamma_{N,n_L}=\frac{2n_L(N-n_L)}{N(N-1)}\longrightarrow2\rho(1-\rho),
\end{equation}
all limit orders give \eqref{eq:exact-Meissner-value}.  Differentiating with respect to $Q=2A$ instead of
$A$ multiplies the curvature by $1/4$ and gives \eqref{eq:pair-phase-value}.
\end{proof}

\subsection{Pair ODLRO and the closing fixed-number gap}

\begin{proposition}[Exact pair-density matrix]
\label{prop:ODLRO}
In the state \eqref{eq:Dicke},
\begin{equation}
 \bra{\Omega_{N,n}}p_x^\dagger p_y\ket{\Omega_{N,n}}
 =\begin{cases}
 n/N,&x=y,\\[0.25em]
 \alpha_{N,n},&x\ne y.
 \end{cases}
 \label{eq:pair-correlator}
\end{equation}
The largest eigenvalue of this matrix is
\begin{equation}
 \lambda_{\max}=\frac{n(N-n+1)}N,
 \label{eq:ODLRO-eigenvalue}
\end{equation}
which is extensive at fixed interior density.  The finite-volume electronic-link Meissner kernel satisfies
\begin{equation}
 j\gamma_{N,n}=2j\bra{\Omega_{N,n}}p_x^\dagger p_y\ket{\Omega_{N,n}}
 \qquad(x\ne y).
 \label{eq:response-ODLRO}
\end{equation}
\end{proposition}

\begin{proof}
The diagonal entry is the probability that a uniformly chosen $n$-subset contains $x$.  For $x\ne y$,
$p_x^\dagger p_y$ maps a subset containing $y$ but not $x$ to another subset with the same uniform
amplitude.  The number of such subsets is $\binom{N-2}{n-1}$, giving $\alpha_{N,n}$.  A matrix with
constant diagonal and constant off-diagonal entries has uniform eigenvalue
$n/N+(N-1)\alpha_{N,n}$, which is \eqref{eq:ODLRO-eigenvalue}.  Equation
\eqref{eq:response-ODLRO} follows from \eqref{eq:alpha-gamma}.
\end{proof}

\begin{proposition}[No hidden uniform gap]
\label{prop:gap-closing}
In the fixed-$n$ paired sector on a periodic nearest-neighbor lattice,
\begin{equation}
 \operatorname{gap}(H_{G_L,n})
 =\frac j4\lambda_1(BB^*)
 =j\sin^2\!\left(\frac\pi L\right)
 =O(L^{-2}),
 \label{eq:gap-closing}
\end{equation}
for $1\le n\le N-1$.  Here $\lambda_1$ is the smallest nonzero graph-Laplacian eigenvalue.
\end{proposition}

\begin{proof}
The upper bound and eigenvalue are realized by the site-function vectors in
\cref{lem:site-functions}.  Equality of the symmetric-exclusion spectral gap with the one-particle
random-walk gap is the Aldous spectral-gap theorem \cite{Caputo2010}.  For the periodic lattice,
$\lambda_1(BB^*)=4\sin^2(\pi/L)$.
\end{proof}

The Meissner theorem is therefore not a uniform-gap corollary.  The gapless modes are longitudinal site
functions.  A transverse bond field has $Ba=0$, so its physical current source vanishes before any
pseudoinverse is applied.

\section{Protected multiblock stiffness floor}
\label{sec:multiblock}

Consider $r$ paired blocks on the same graph.  Let the full positive-square target map be stacked as
\begin{equation}
 D[A]=D_{\rm sw}[A]\oplus D_{\rm r}[A],
 \label{eq:stacked-map}
\end{equation}
where $D_{\rm sw}$ contains a complete swap row for every graph edge in each block, with coupling $j_b$,
and every row of $D_{\rm r}$ annihilates the product-Dicke zero manifold at $A=0$.  Let
$S_{\rm sw}$ and $S_{\rm r}$ be the corresponding Peierls source maps.

\begin{theorem}[Source-kernel backbone protection]
\label{thm:backbone}
For every divergence-free profile $Ba=0$ and every $\zeta>0$,
\begin{align}
 \cK(\I\zeta;a)&\succeq Q_{\rm sw}[a],
 \label{eq:backbone-dynamic}\\
 \cC[a]&\succeq Q_{\rm sw}[a],
 \label{eq:backbone-static}
\end{align}
where, in the product-composition basis $\ket{\bm n}$,
\begin{equation}
 Q_{\rm sw}[a]
 =\norm a^2\sum_{\bm n}
 \left(\sum_{b=1}^rj_b\gamma_{N,n_b}\right)
 \ket{\bm n}\bra{\bm n}.
 \label{eq:backbone-floor}
\end{equation}
Thus additional positive-square multiplicative, remote-control, or interaction rows cannot cancel the
transverse stiffness carried by the swap backbone.
\end{theorem}

\begin{proof}
For $Ba=0$, the single-block theorem gives
$D_{\rm sw}^\dagger S_{\rm sw}=0$.  Because the target labels are a direct sum,
$s_0:=(S_{\rm sw},0)\in\ker D^\dagger=\ker\cL$ for the complete map.  Hence
$f_\zeta(\cL)s_0=s_0$ for
$f_\zeta(\cL)=\zeta^2(\cL^2+\zeta^2)^{-1}$.  Although $\cL=DD^\dagger/2$ need not preserve the
swap/residual block decomposition, self-adjointness now gives, for
$s_{\rm r}:=(0,S_{\rm r})$,
\begin{equation}
 \ip{s_0}{f_\zeta(\cL)s_{\rm r}}
 =\ip{f_\zeta(\cL)s_0}{s_{\rm r}}
 =\ip{s_0}{s_{\rm r}}=0.
\end{equation}
Thus it is the kernel eigenvector property, followed by direct-sum orthogonality, that removes the cross
terms.  Applying \cref{thm:dynamical-LS} gives the direct positive swap Gram matrix plus a positive
residual term.  The static statement follows identically with $P_{\ker D^\dagger}$ in place of
$f_\zeta(\cL)$.  Different blocks preserve separate block charges and occupy orthogonal row labels,
yielding \eqref{eq:backbone-floor}.
\end{proof}

\begin{corollary}[Selected interior compositions are Meissner active]
\label{cor:selected-composition}
On a periodic $d\ge2$ lattice, choose a product composition with $n_{b,L}/N\to\rho_b$.  Then
\begin{equation}
 D_{\rm M}\ge2\sum_{b=1}^rj_b\rho_b(1-\rho_b),
 \label{eq:multiblock-floor}
\end{equation}
and the same lower bound holds for the Abelian zero-frequency kernel.  It is strictly positive whenever at
least one block is neither empty nor saturated.
\end{corollary}

\begin{remark}[Ground-manifold selection]
In an unresolved composition manifold, the lowest eigenvalue of the full curvature operator need not
equal its diagonal expectation in a chosen product composition.  A scalar physical response therefore
requires a branch selector, a prepared coherent state, or a density matrix.  The operator inequalities
\eqref{eq:backbone-dynamic}--\eqref{eq:backbone-static} remain valid without choosing one.
\end{remark}

\section{Independent failure mechanisms}
\label{sec:counterexamples}

\subsection{Temporal failure from a soft positive target mode}

\begin{proposition}[Positive-square order-of-limits defect]
\label{prop:soft-target-counterexample}
Let $V_L\to\infty$, $\Delta_L\downarrow0$, and let the state space be spanned by
$\ket g,\ket e$ and the target space by $\ket r,\ket z$.  Define
\begin{equation}
 D_L(a)
 =\sqrt{2\Delta_L}\ket r\bra e
 +a\left(\sqrt{2\alpha V_L}\ket r+\sqrt{c_0V_L}\ket z\right)\bra g,
 \qquad
 \alpha>0,
 \quad c_0\ge0,
 \label{eq:soft-target-D}
\end{equation}
and $H_L(a)=D_L(a)^\dagger D_L(a)/2$.  At $a=0$, $\ket g$ is the unique zero state and $\ket e$ has
energy $\Delta_L$.  Nevertheless,
\begin{align}
 \frac1{V_L}\cC_L&=c_0,
 \label{eq:soft-static}\\
 \frac1{V_L}\cK_L(\I\zeta)
 &=c_0+2\alpha\frac{\zeta^2}{\Delta_L^2+\zeta^2}.
 \label{eq:soft-dynamic}
\end{align}
Hence
\begin{equation}
 \lim_{L\to\infty}\frac{\cC_L}{V_L}=c_0,
 \qquad
 \lim_{\zeta\downarrow0}\lim_{L\to\infty}
 \frac{\cK_L(\I\zeta)}{V_L}=c_0+2\alpha.
 \label{eq:soft-noncommute}
\end{equation}
\end{proposition}

\begin{proof}
At zero twist,
$\cL_L=\Delta_L\ket r\bra r$ and the source is the second term in
\eqref{eq:soft-target-D} acting on $\ket g$.  The static projector onto $\ker\cL_L=\Span\{\ket z\}$
gives \eqref{eq:soft-static}; \cref{thm:dynamical-LS} gives \eqref{eq:soft-dynamic}.  Taking the limits
proves \eqref{eq:soft-noncommute}.  The soft mode carries source weight $2\alpha V_L$ at target energy
$\Delta_L$, so \eqref{eq:source-tightness} fails exactly.
\end{proof}

\subsection{Spatial failure: Drude without Meissner}

\begin{proposition}[Clean normal metal]
\label{prop:free-gas}
Consider continuum spinless fermions on a periodic box,
\begin{equation}
 H_L(A)=\sum_{a=1}^N\frac{(p_a-A)^2}{2m},
 \label{eq:free-gas}
\end{equation}
with a closed-shell ground state and density $n=N/V_L$.  The finite-volume Kohn curvature and uniform
Drude kernel are $n/m$.  The static transverse kernel instead vanishes in the thermodynamic limit:
\begin{equation}
 D_{\rm M}=0,
 \qquad
 \kappa_{\rm K}=D_{\rm A}=\frac nm.
 \label{eq:free-gas-result}
\end{equation}
\end{proposition}

\begin{proof}
Along a fixed zero-temperature closed-shell adiabatic branch, differentiating \eqref{eq:free-gas} twice
gives $n/m$.  To evaluate the transverse thermodynamic limit without distributional shorthand, first use
the smooth Fermi occupation $n_\beta(\epsilon)=[1+\e^{\beta(\epsilon-\mu)}]^{-1}$ at finite
$\beta$.  For $q$ along $x$ and polarization along $y$,
\begin{equation}
 \Pi_{yy}^{(\beta)}(q,0)
 =\int\frac{\dd^dk}{(2\pi)^d}
 \left(\frac{k_y}{m}\right)^2
 \frac{n_\beta(\epsilon_{k+q/2})-n_\beta(\epsilon_{k-q/2})}
 {\epsilon_{k+q/2}-\epsilon_{k-q/2}}.
\end{equation}
At fixed $\beta$, dominated convergence as $q\to0$ gives the derivative
$n_\beta'(\epsilon_k)$.  The identity
\begin{equation}
 \partial_{k_y}\!\left(\frac{k_y}{m}n_\beta(\epsilon_k)\right)
 =\frac1m n_\beta(\epsilon_k)+\frac{k_y^2}{m^2}n_\beta'(\epsilon_k)
\end{equation}
and integration by parts yield
$\lim_{q\to0}\Pi_{yy}^{(\beta)}(q,0)=-n_\beta/m$, where $n_\beta$ is the particle density at that
regularization.  Taking $\beta\to\infty$ at fixed density gives $-n/m$, equivalently the usual
zero-temperature Fermi-surface distribution.  It cancels the diamagnetic term.  The finite-$\beta$
numerical check in the reproducibility suite tests this smooth regularization; the proposition is its
zero-temperature limit.  The spatial continuity hypothesis of \cref{thm:general-completion} therefore
fails.
\end{proof}

These examples show why both theorem layers are necessary.  The first violates temporal tightness in a
frustration-free positive-square family.  The second satisfies the uniform electric response but violates
the transverse limit.

\section{Theorem boundary and relation to microscopic transfer}
\label{sec:boundary}

The exact parent has now crossed the conceptual boundary between static Peierls curvature and a genuine
zero-frequency transverse matter response.  The proof uses five pieces that can be audited separately:

\begin{center}
\begin{tabular}{@{}>{\raggedright\arraybackslash}p{0.27\linewidth}>{\raggedright\arraybackslash}p{0.31\linewidth}>{\raggedright\arraybackslash}p{0.32\linewidth}@{}}
\toprule
Question & Sufficient certificate & Failure signature\\
\midrule
Finite-volume branch
& isolated ground cluster and $PJP=0$
& linear splitting, cusp, or unresolved selector\\
Complete gauge coupling
& bondwise covariance and both current/stress derivatives
& partial twist or missing diamagnetic term\\
Temporal uniformity
& current $H^{-1}$ tightness, equivalently source tightness
& finite inverse-energy weight in soft modes\\
Transverse completion
& static $q\to0$ continuity or exact source-kernel identity
& Drude response with transverse cancellation\\
Exact swap route
& $Ba=0$ and $S[a]\subset\ker D^\dagger$
& source outside the target kernel\\
\bottomrule
\end{tabular}
\end{center}

The finite-coupling microscopic question is different.  A local Schrieffer--Wolff low block can reproduce
the parent at leading order while carrying a quasi-local higher-order remainder.  Because the parent gap
closes as $L^{-2}$, a small local operator norm does not by itself imply a small response error: the
remainder can move the selected state or couple a transverse source to a soft longitudinal mode.  The
required transfer theorem must control both state and current sources in a relative $H^{-1}$ topology,
uniformly in the transverse infrared window.  That analysis is intentionally separated from the exact
Meissner theorem proved here.

\section{Conclusion}

Static Peierls curvature becomes a dynamical statement only after an infrared uniformity theorem.  The
sharp temporal condition is not that the complete many-body gap stay open, but that the physical current
carry no residual inverse-energy weight into asymptotically soft excitations.  In a positive-square parent,
this is exactly a statement about the projection of the Peierls source onto soft positive target singular
modes.

A Meissner statement additionally requires the transverse spatial limit.  The graph swap parent solves
that problem at finite size.  Its current source is the graph divergence of the bond field; least-squares
relaxation is Hodge projection; and every divergence-free field lies exactly in the target kernel.  The
transverse kernel is therefore independent of frequency and wave vector before any thermodynamic limit
is taken.

The resulting periodic-lattice model has
\begin{equation}
 \kappa_{\rm K}=D_{\rm A}=D_{\rm M}=2j\rho(1-\rho)>0
\end{equation}
at every interior filling, exact pair ODLRO, and a fixed-number gap closing as $L^{-2}$.  It is thus an
explicit gapless frustration-free parent with a rigorous zero-frequency Meissner/superfluid weight.  The
swap backbone also protects a positive transverse floor in multiblock extensions.  The next independent
problem is not conceptual identification but response-stable transfer to a specified finite-coupling
microscopic realization.

\appendix

\section{Finite-volume perturbation formula}
\label{app:perturbation}

For completeness, write
\begin{equation}
 H_L(a)=H_L(0)+aJ_L+\frac{a^2}{2}T_L+O(a^3).
\end{equation}
With respect to $P_L\oplus Q_L$ and under $P_LJ_LP_L=0$, the Feshbach map at energies $O(a^2)$ is
\begin{equation}
 \frac{a^2}{2}P_LT_LP_L
 -a^2P_LJ_LQ_L(Q_L\widehat H_LQ_L)^{-1}Q_LJ_LP_L
 +o(a^2).
\end{equation}
Factoring $a^2/2$ gives \eqref{eq:Kohn-operator}.  Thus the analytic second-order branches emerging
from a degenerate ground cluster are the eigenvalues of $\cC_L$.  The lowest unresolved ground-energy
curvature is $\lambda_{\min}(\cC_L)$, while a prepared state or density matrix evaluates the corresponding
expectation.  These are different physical selection rules when the cluster is degenerate.

\section{Graph combinatorics}
\label{app:graph-combinatorics}

The subset average used in \cref{lem:site-functions} follows from
\begin{align}
 \mathbb P(x\in S)&=\frac nN,\\
 \mathbb P(x,y\in S)&=\frac{n(n-1)}{N(N-1)},\qquad x\ne y.
\end{align}
For zero-mean $f,g$,
\begin{align}
 \mathbb E\left[\sum_{x\in S}f_x\sum_{y\in S}g_y\right]
 &=\frac nN\sum_xf_xg_x
 +\frac{n(n-1)}{N(N-1)}\sum_{x\ne y}f_xg_y\\
 &=\frac{n(N-n)}{N(N-1)}\sum_xf_xg_x.
\end{align}
The same counting gives the probability that exactly one endpoint of an edge is occupied:
\begin{equation}
 2\frac nN\frac{N-n}{N-1}=\gamma_{N,n}.
\end{equation}

\section{Normalization conventions}
\label{app:normalization}

The electronic-link convention uses the Peierls phase of a single electron and the response density
$V^{-1}\partial_A^2E$.  A pair transported across a bond acquires phase $Q=2A$.  Therefore
\begin{equation}
 \frac1V\frac{\partial^2E}{\partial Q^2}
 =\frac1{4V}\frac{\partial^2E}{\partial A^2}.
\end{equation}
This accounts for the factor of four between \eqref{eq:exact-Meissner-value} and
\eqref{eq:pair-phase-value}.  Electric charge factors, optical $\pi$ factors, and lattice-spacing tensors can
be restored uniformly after choosing a material convention.

\section{Reproducibility}
\label{app:reproducibility}

Three deterministic scripts accompany the paper.

The script \texttt{verify\_dynamical\_response.py} checks:
\begin{enumerate}[leftmargin=1.6em]
\item random finite-dimensional positive-square families against the state-space Kubo formula and the
target-space identity \eqref{eq:dynamical-LS};
\item the quantitative bounds \eqref{eq:lower-tail-bound}--\eqref{eq:upper-tail-bound} for random
positive matrix-valued spectral weights;
\item the soft-target counterexample \eqref{eq:soft-static}--\eqref{eq:soft-noncommute};
\item numerical convergence of the free-gas transverse cancellation.
\end{enumerate}

The script \texttt{verify\_graph\_hodge\_meissner.py} checks:
\begin{enumerate}[leftmargin=1.6em]
\item the unique Dicke/AGP zero state and its pair-density matrix;
\item the site-function identities \eqref{eq:site-inner}--\eqref{eq:site-Laplacian};
\item bondwise gauge covariance;
\item the current and stress identities \eqref{eq:current-source}--\eqref{eq:stress-expectation};
\item the static Hodge formula and finite-frequency graph filter;
\item exact cancellation for gradients and exact survival for lattice-transverse profiles;
\item the $L^{-2}$ gap formula, including a multipair half-filled rectangular torus, and an operator-level multiblock floor with nontrivial residual rows.
\end{enumerate}

The script \texttt{independent\_fock\_check.py} checks the fermionic layer before seniority-zero
compression.  It constructs Jordan--Wigner Fock matrices on a path, triangle, and square and verifies:
\begin{enumerate}[leftmargin=1.6em]
\item the complete site swaps and their Peierls lifts are Hermitian involutions;
\item the even-sector kernels have the claimed fermionic Dicke signs and odd sectors have no zero modes;
\item full-Fock gauge covariance and the charge-two seniority-zero compression;
\item current and stress derivatives against finite differences; and
\item the full-Fock static Hodge formula and finite-frequency transverse identity.
\end{enumerate}
The retained outputs and SHA-256 hashes are included in the consolidated archive.

\begingroup\raggedright
\begin{thebibliography}{99}

\bibitem{Kohn1964}
W.~Kohn,
``Theory of the Insulating State,''
\emph{Phys. Rev.} \textbf{133}, A171--A181 (1964),
doi:10.1103/PhysRev.133.A171.

\bibitem{Scalapino1992}
D.~J. Scalapino, S.~R. White, and S.~C. Zhang,
``Superfluid density and the Drude weight of the Hubbard model,''
\emph{Phys. Rev. Lett.} \textbf{68}, 2830--2833 (1992),
doi:10.1103/PhysRevLett.68.2830.

\bibitem{Scalapino1993}
D.~J. Scalapino, S.~R. White, and S.~C. Zhang,
``Insulator, metal, or superconductor: The criteria,''
\emph{Phys. Rev. B} \textbf{47}, 7995--8007 (1993),
doi:10.1103/PhysRevB.47.7995.

\bibitem{Resta2018}
R.~Resta,
``Drude weight and superconducting weight,''
\emph{J. Phys.: Condens. Matter} \textbf{30}, 414001 (2018),
doi:10.1088/1361-648X/aade19.

\bibitem{Watanabe2020}
H.~Watanabe and M.~Oshikawa,
``Generalized f-sum rules and Kohn formulas on nonlinear conductivities,''
\emph{Phys. Rev. B} \textbf{102}, 165137 (2020),
doi:10.1103/PhysRevB.102.165137.

\bibitem{Takasan2023}
K.~Takasan, M.~Oshikawa, and H.~Watanabe,
``Drude weights in one-dimensional systems with a single defect,''
\emph{Phys. Rev. B} \textbf{107}, 075141 (2023),
doi:10.1103/PhysRevB.107.075141.

\bibitem{Tada2016}
Y.~Tada and T.~Koma,
``Two No-Go Theorems on Superconductivity,''
\emph{J. Stat. Phys.} \textbf{165}, 455--470 (2016),
doi:10.1007/s10955-016-1629-2.

\bibitem{Peotta2015}
S.~Peotta and P.~T{\"o}rm{\"a},
``Superfluidity in topologically nontrivial flat bands,''
\emph{Nat. Commun.} \textbf{6}, 8944 (2015),
doi:10.1038/ncomms9944.

\bibitem{Torma2022}
P.~T{\"o}rm{\"a}, S.~Peotta, and B.~A. Bernevig,
``Superconductivity, superfluidity and quantum geometry in twisted multilayer systems,''
\emph{Nat. Rev. Phys.} \textbf{4}, 528--542 (2022),
doi:10.1038/s42254-022-00466-y.

\bibitem{Sewell1990}
G.~L. Sewell,
``Off-diagonal long range order and the Meissner effect,''
\emph{J. Stat. Phys.} \textbf{61}, 415--422 (1990),
doi:10.1007/BF01013973.

\bibitem{Nieh1995}
H.~T. Nieh, G.~Su, and B.-H. Zhao,
``Off-diagonal long-range order: Meissner effect and flux quantization,''
\emph{Phys. Rev. B} \textbf{51}, 3760--3764 (1995),
doi:10.1103/PhysRevB.51.3760.

\bibitem{Matsubara1956}
T.~Matsubara and H.~Matsuda,
``A Lattice Model of Liquid Helium, I,''
\emph{Prog. Theor. Phys.} \textbf{16}, 569--582 (1956),
doi:10.1143/PTP.16.569.

\bibitem{Yang1962}
C.~N. Yang,
``Concept of Off-Diagonal Long-Range Order and the Quantum Phases of Liquid He and of Superconductors,''
\emph{Rev. Mod. Phys.} \textbf{34}, 694--704 (1962),
doi:10.1103/RevModPhys.34.694.

\bibitem{Yang1989}
C.~N. Yang,
``$\eta$ Pairing and Off-Diagonal Long-Range Order in a Hubbard Model,''
\emph{Phys. Rev. Lett.} \textbf{63}, 2144--2147 (1989),
doi:10.1103/PhysRevLett.63.2144.

\bibitem{Caputo2010}
P.~Caputo, T.~M. Liggett, and T.~Richthammer,
``Proof of Aldous' Spectral Gap Conjecture,''
\emph{J. Amer. Math. Soc.} \textbf{23}, 831--851 (2010),
doi:10.1090/S0894-0347-10-00659-4.

\bibitem{Han2024}
Z.~Han, J.~Herzog-Arbeitman, B.~A. Bernevig, and S.~A. Kivelson,
``Quantum Geometric Nesting and Solvable Model Flat-Band Systems,''
\emph{Phys. Rev. X} \textbf{14}, 041004 (2024),
doi:10.1103/PhysRevX.14.041004.

\bibitem{Gao2026}
Q.~Gao, Z.~Han, and E.~Khalaf,
``Bootstrapping Flat-band Superconductors: Rigorous Lower Bounds on Superfluid Stiffness,''
\emph{Phys. Rev. Lett.} \textbf{136}, 076503 (2026),
doi:10.1103/gw85-5r92.

\bibitem{SledgianowskiSingle}
N.~Sledgianowski,
``Exact Flat-Band Stiffness from Pair Mobility in a QGN Model,''
manuscript (2026).

\bibitem{SledgianowskiMulti}
N.~Sledgianowski,
``Multi-Block Reduction and Irreducibility Laws for QGN Superconductors,''
manuscript (2026).

\bibitem{SledgianowskiNoGo}
N.~Sledgianowski,
``Peierls-Transport Obstruction and the Minimal Many-Body Escape,''
manuscript (2026).

\bibitem{SledgianowskiMicro}
N.~Sledgianowski,
``Microscopic Origins of Multiplicative Null Interactions in Quantum-Geometric-Nesting Superconductors,''
manuscript (2026).

\end{thebibliography}
\endgroup

\end{document}
